\section{The corrected reachability predicate}

This chapter explains the corrected reachability predicate
$\ResetWindow(j,k_0,t,\delta,s)$ field by field, and why each field is
necessary.

\subsection{The fields}

The predicate $\ResetWindow(j,k_0,t,\delta,s)$ asserts the existence
of states $r$ and $r_j$ such that:
\begin{enumerate}[leftmargin=*]
\item $s5^{k_0}=r+1$;
\item $\GeneralOrbit(r)$;
\item $\FullOrbit(r_j)$;
\item $\mathrm{ResetHeadEq}(s,j,k_0,t,\delta,r_j)$;
\item $s<5^{j-1-k_0}$;
\item $s$ is odd;
\item $5\nmid s$.
\end{enumerate}
The first four items are reachability data; the last three are size
and parity conditions.

\subsection{Why the same indices}

The predicate fixes the same $j,k_0,t$ in every field.  The reset
equation, the size bound, and the two reachability predicates all use
the same depth and weight.  This prevents a witness from combining a
reset equation at one depth with a size bound at another.

\subsection{A weaker predicate would fail}

A predicate that only fixed the difference $j-k_0-1$ would allow a
reset equation with different absolute indices than the reachability
witness.  The corrected predicate is written with the absolute
indices precisely to block that gap.

\subsection{The arithmetic witness}

The witness
\begin{equation*}
  j=36,\quad k_0=0,\quad t=2,\quad \delta=3,\quad
  s=2^{83}-3\cdot5^{35}
\end{equation*}
satisfies the size and parity conditions:
\begin{itemize}[leftmargin=*]
\item $s<5^{35}$ because $2^{81}<5^{35}$;
\item $s$ is odd;
\item $s\equiv3\pmod5$, so $5\nmid s$.
\end{itemize}
It also satisfies the valuation identity
\begin{equation*}
  5s+3\cdot5^{36}=5\cdot2^{83},
\end{equation*}
which is why it breaks the arithmetic-only bound.

\subsection{What the witness lacks}

The witness supplies no $r$ with
\begin{equation*}
  s5^{0}=r+1,
\end{equation*}
no general-orbit reachability $\GeneralOrbit(r)$, no full-orbit reset
head $\FullOrbit(r_j)$, and no reset equation
\begin{equation*}
  \mathrm{ResetHeadEq}(s,36,0,2,3,r_j).
\end{equation*}
It is arithmetic-only, so it is not a counterexample to the corrected
bound.

\subsection{Statement shape}

The repository definition is
\texttt{S6Audit.ResetWindowReachability}.  The paper writes it
schematically as an existence predicate over $r$ and $r_j$ with the
seven fields listed above; the precise record structure lives in the
repository and may use different field names.  The mathematical
content is the seven fields, not the layout of the record.

\subsection{Status}

\begin{center}
\small
\begin{tabular}{@{}p{0.56\textwidth}p{0.36\textwidth}@{}}
\toprule
Item & Status \\
\midrule
predicate definition & compiled \\
arithmetic witness & verified exactly \\
corrected bound using the predicate & compiled (conditional) \\
failure-window construction & pending \\
\bottomrule
\end{tabular}
\end{center}
